\begin{satz}{Satz von Gluskin}{SatzVonGluskin}
Seien $m \in \N$, $\eps \in\!\ ]0, 1[$, $N\coloneqq  \ceil{(1-\eps)m}$ und
\[c \coloneqq  2\sqrt \pi (\sqrt 2 +1) \sqrt{2 + \frac 4{\log(4)}}
\sqrt{\frac{\log\left(\pi(\sqrt 2+1)^2\left(1+\frac 2{\log(4)}\right) \right)}
{\log(2)} + 1}\text.\]
Dann existiert $F \le \ell_1^m$ mit $\dim F = m - N$ und
\[d_{BM}(F,\ell_2^{m-N})<
c\sqrt{\frac 1{1-\eps}\log\left(\frac 2{1-\eps}\right)}\text.\]
\end{satz}
\begin{proof}
Seien $\rho \coloneqq  (\sqrt 2 +1)^2\left(2+\frac 4{\log(4)} \right)$,
$\eta \coloneqq  \frac 2 \pi \frac 1 \rho (1-\eps)$ und
$t_\eta \coloneqq  \frac 4 \eta \log\left(\frac 2 \eta \right)$.\\
Dann ist $c^2=4\pi\rho \left(\frac{\log\left(\frac{\pi\rho}2\right)}{\log(2)}+1\right)$ und damit:
\begin{align}
&2t_\eta = \frac{4\pi\rho}{1-\eps}\log\left(\frac{\pi\rho}{1-\eps} \right)
= \frac{4\pi\rho}{1-\eps}\log\left(\frac{\pi\rho}{2}\frac{2}{1-\eps} \right)
\notag\\
=& \frac{4\pi\rho}{1-\eps}\left(\log\left(\frac{\pi\rho}2 \right) 
 + \log\left( \frac 2 {1-\eps}\right) \right)
 \overset{\eps>-1}=
 \frac{4\pi\rho}{1-\eps}\left(\frac{\log\left(\frac{\pi\rho}2
 \right)}{\log\left( \frac 2 {1-\eps}\right)}+1 \right)\log\left(\frac 2
 {1-\eps} \right)\notag\\
 =&4\pi\rho\left(\frac{\log\left(\frac{\pi\rho}2 \right)}
 {\log\left( \frac 2 {1-\eps}\right)}+1 \right)
 \frac 1{1-\eps}\log\left(\frac 2 {1-\eps} \right)
 \overset{\frac 2{1-\eps}>2}<
 c^2
 \frac 1{1-\eps}\log\bigg(\frac 2 {1-\eps} \bigg). \label{svgeq:1}
\end{align}
Falls $N=m$ gilt, ist die Aussage trivial, sodass wir $N<m$ annehmen können.\\
\underline{1. Fall:} $m \le 2 t_\eta$.\\
Sei $F \coloneqq  \left\langle \menge{e_i^m}{i \in \haken {m-N}}\right\rangle$.\\
Offenbar ist dann $(F,\restrict{\norm\cdot_1}{F})$ isometrisch isomorph zu
$\ell_1^{m-N}$ und damit nach~\refclaim{EigenschaftenDesBanachMazurAbstands}
{EigenschaftenDesBanachMazurAbstands2}:
$d_{BM}((F,\restrict{\norm\cdot_1}{F}), \ell_1^{m-N}) = 1$.\\
Weiter gilt nach
\refclaim{AequivalenzkonstantenFuerPNormen}{AequivalenzkonstantenFuerPNormen1}
$\norm\cdot_{2, m-N} \le\norm\cdot_{1, m-N} \le \sqrt{m-N}\norm\cdot_{2, m-N}$
und somit nach~\refclaim{EigenschaftenDesBanachMazurAbstands}
{EigenschaftenDesBanachMazurAbstands3}: $d_{BM}(\ell_1^{m-N}, \ell_2^{m-N}) \le
\sqrt{m-N}$.\\
Somit gilt
\begin{align*}
d_{BM}\left((F,\restrict{\norm\cdot_1}F), \ell_2^{m-N}\right)
&\overset{\refclaim{EigenschaftenDesBanachMazurAbstands}
{EigenschaftenDesBanachMazurAbstands1}}\le
d_{BM}((F,\restrict{\norm\cdot_1}{F}), \ell_1^{m-N})d_{BM}
(\ell_1^{m-N}, \ell_2^{m-N})\\
& \le \sqrt{m-N} \le \sqrt{m - (1-\eps)m} = \sqrt{\eps m} < \sqrt{m}
\le \sqrt{2t_\eta}\\
&\overset{\eqref{svgeq:1}}\le
c \sqrt{\frac 1 {1-\eps}\log\left(\frac 2{1-\eps} \right)}
\text.
\end{align*}
Damit ist also $F$ ein Unterraum wie behauptet.\\
\underline{2. Fall:} $m > 2t_\eta$.\\
Dann ist $1<m\frac1{2t_\eta}=m\left(\frac1{t_\eta}-\frac1{2t_\eta}\right)$ und
damit $\N \cap \left[\frac 1 2 \frac m {t_\eta}, \frac m {t_\eta}\right]
\neq \emptyset$.\\
Sei $k\in\N\cap\left[\frac12\frac m{t_\eta},\frac m{t_\eta}\right]$.\\
Offenbar ist $\eta < 1$ und damit $t_\eta > 4 \log(2) = \log(16) > 2$.\\
Also ist $k \le \frac m {t_\eta} < \frac m 2$.\\
Seien $(E,\norm\cdot_E)\coloneqq \ell_2^N$, $T\coloneqq S^{m-1}\cap\kabg{\ell_1^m} 0{\sqrt k}$,
$x \coloneqq  (e_i)_{i = 1}^N$,
\[\bigg((\Omega, \mfa, \Prim), (\gamma,\delta,\widetilde g),(N, m, mN)\bigg)
\text{ ein }(N(0,1), 3)\text{--Tupel zu }(3,(N,m,mN))\] (existiert nach~\ref{ExistenzNormalNTupel}) und $(Z_{\gamma, x}, M_{\delta,T})$ das
Auswertungspaar zu $(\gamma, \delta, x, T)$.\\
Dann gilt $\bprod \cdot x_N = \id_{\R^N}$ und damit $x \neq 0$ und
$\norm\cdot_E\circ\bprod\cdot x _N=\norm\cdot_{2,N}$. \hfill $(a)$\\
Offenbar gilt $T \neq \emptyset$ (wegen $k \ge 1$ ist $e_1^m \in T$),
$T \subseteq S^{m-1}$, $T=-T$. \hfill $(b)$\\
Ferner ist $T$ als Schnitt $1$--symmetrischer Mengen selbst $1$--symmetrisch
und damit gilt ($\mu_{i,n}$ für alle $i \in \haken m$ gemäß~\ref{ktesMaximum})
\begin{align}
\forall\ a \in \R^m : \sup_{t \in T} \sprod t a_m = \sup_{t \in T} \sprod
t {((\mu_{i,n}\circ \pfeil {\abs\cdot} m)(a))_{i=1}^m}_m \text.\label{svgeq:2}
\end{align}
Sei $\omega \in \Omega$. Dann gilt für alle $s \in T$:
\begin{align}
 \sum_{i=k+1}^m s_i\delta_i^\ast(\omega) 
&\le \sum_{i=k+1}^m \abs{s_i}\delta_i^\ast(\omega) 
\overset{\refclaim{EigenschaftenDesKtenMaximums}
{EigenschaftenDesKtenMaximums1}}\le
\delta_{k+1}^\ast(\omega) \sum_{i=k+1}^m \abs{s_i} \notag\\
&\overset{\refclaim{EigenschaftenDesKtenMaximums}
{EigenschaftenDesKtenMaximums1}}\le \delta_k^\ast(\omega) \norm s _{1,m}
\le \sqrt k \cdot \delta_k^\ast(\omega) \text.\label{svgeq:3}
\end{align}
Weiter gilt wegen $k \le \frac m {t_\eta}$ auch $\frac m k \ge t_\eta$.\\
Da $\eta \in\!\  ]0,1[$ gilt nach Setzung von
$t_\eta$ gemäß~\refclaim{Numeriklemma}{Numeriklemma4}:
\begin{align}
\log\left(\frac{2m}k \right) < \eta \frac m k\text. \label{svgeq:4}
\end{align}
Somit gilt:
\begin{align*}
&\,\E\left(M_{\delta,T}\right)=\E\left(\sup_{t \in T} \bprod t \delta_m \right)
\overset{\eqref{svgeq:2}}=
\E\left(\sup_{t\in T}\bprod t{\delta^\ast}_m \right)\\
= &\, \E\left(\sup_{t \in T}\left(\left(\sum_{i=1}^k t_i\delta_i^\ast\right) 
+ \left(\sum_{i=k+1}^m t_i\delta_i^\ast \right) \right) \right)\\
\le &\, \E\left(\sup_{t \in T}\sum_{i=1}^k t_i\delta_i^\ast \right) 
+ \E\left(\sup_{t \in T} \sum_{i=k+1}^m t_i\delta_i^\ast\right) \\
\overset{\eqref{svgeq:3}}\le &\,
\E\left(\sup_{t \in T}\sum_{i=1}^k t_i\delta_i^\ast \right)
+ \sqrt k \cdot \E\left(\delta_k^\ast \right)\\
\le &\, \E\left(\sup_{t \in T}\norm{t}_{2,m}\left(
\norm\cdot_{2,k} \circ \left(\delta_i^\ast \right)_{i=1}^k \right) \right)
+ \sqrt k \cdot \E\left(\delta_k^\ast \right)\\
\overset{T \subseteq S^{m-1}}=
&\, \E\left(\norm\cdot_{2,k} \circ \left(\delta_i^\ast \right)_{i=1}^k \right)
+ \sqrt k \cdot \E\left(\delta_k^\ast \right)\\
\overset{\substack{\refclaim
{ObereAbschaetzungDesErwartungswertsDerKtenAbsolutenN01Ordnungsstatistik}
{ObereAbschaetzungDesErwartungswertsDerKtenAbsolutenN01Ordnungsstatistik5},\\
\refclaim
{ObereAbschaetzungDesErwartungswertsDerKtenAbsolutenN01Ordnungsstatistik}
{ObereAbschaetzungDesErwartungswertsDerKtenAbsolutenN01Ordnungsstatistik7}}}
\le&\,\sqrt 2\sqrt{2+ \frac 4 {\log(4)}} \sqrt{k \log\left(\frac {2m}k\right)}
+\sqrt k \cdot\sqrt{2 + \frac 4 {\log(4)}}\sqrt{\log\left(\frac {2m}k\right)}\\
=&\, \left(\sqrt{2} + 1\right)\sqrt{2+ \frac 4 {\log(4)}}
\sqrt{k \log\left(\frac {2m}k\right)}\\
\overset{\eqref{svgeq:4}}<
&\, \left(\sqrt{2} + 1\right)\sqrt{2+ \frac 4 {\log(4)}}
\sqrt{k \eta \frac mk}\\
=&\, \left(\sqrt{2} + 1\right)\sqrt{2+ \frac 4 {\log(4)}}
\sqrt{\eta m} = 
\sqrt{\frac 2 \pi} \sqrt{(1-\eps)m} \le \sqrt{\frac 2 \pi}
\sqrt N \overset{\refclaim{AbschaetzungEinesGammafunktionsquotienten}
{AbschaetzungEinesGammafunktionsquotienten4}}\le a_N\\
\overset{~\ref{ErwartungswertDer2NormEinesStandardnormalverteiltenProzesses}}
=&\, \E\left(\norm\cdot_{2, N} \circ \gamma \right)
= \E\left(\norm\cdot_{2, N}\circ \bprod \cdot x_N \circ \gamma \right)
= \E\left(\norm\cdot_E \circ \bprod \cdot x_N \circ \gamma \right) \tag{$c$}
\end{align*}
Also existiert nach~\ref{ExistenzEingeschraenkterFastIsometrien} ein
$A \in L(\ell_2^m, E)$ mit $\inf\menge{\norm{At}_{2, N}}{t \in T} >0$.\\
Wegen $\dim \Kern A = m - \dim \Bild A \ge m - N$ existiert ein $F \le \Kern A$
mit $\dim F = m-N$.\\
Sei $w \in F\setminus \meng 0$.\\
Dann ist $w \in \R^m$ und damit nach
\refclaim{AequivalenzkonstantenFuerPNormen}{AequivalenzkonstantenFuerPNormen1}:
$\norm w _{1,m} \le \sqrt m \norm w _{2,m}$.\\
Wegen $w \neq 0$ ist $\frac 1 {\norm w_{2,m}}w\in S^{m-1}$, also wegen $Aw = 0$
bereits $w \in \R^m \setminus T$.\\
Somit gilt nach Setzung von $T$:
\[\norma{\frac 1 {\norm w_{2,m}}w}_{1,m}>\sqrt k \text{, }\,\text{ also }\,
 \norm w _{1, m} > \sqrt k \norm w _{2,m}\text.\]
Daher gilt: $\forall\, v \in F: \sqrt k \norm v_{2,m} \le \norm v_{1,m}
\le \sqrt m \norm v _{2, m}$.\\
Also gilt
\begin{align}
d_{BM}\left((F,\restrict{\norm\cdot_{1,m}}F), 
(F, \restrict{\norm\cdot_{2,m}}F) \right)
\overset{\refclaim{EigenschaftenDesBanachMazurAbstands}
{EigenschaftenDesBanachMazurAbstands3}} \le \sqrt{\frac m k} \text.
\label{svgeq:5}
\end{align}
Da $F \le \R^m$, ist $\left(F, 
\restrict{\sprod\cdot{\cdot\cdot}_m}{F \times F}\right)$ ein
Hilbertraum und damit sind $\left(F,\restrict{\norm\cdot_{2,m}}F\right)$
und $\ell_2^{m-N}$ wegen gleicher Dimension isometrisch isomorph.\\
Wegen $k \ge \frac 1 2 \frac m {t_\eta}$ gilt schließlich
\begin{align*}
&d_{BM}\left(\left(F,\restrict{\norm\cdot_{1,m}}F\right),\ell_2^{m-N} \right)\\
\overset{\refclaim{EigenschaftenDesBanachMazurAbstands}
{EigenschaftenDesBanachMazurAbstands1}}
\le &\, d_{BM}\left(\left(F,\restrict{\norm\cdot_{1,m}}F\right),
                    \left(F,\restrict{\norm\cdot_{2,m}}F\right) \right)
    d_{BM}\left(\left(F,\restrict{\norm\cdot_{2,m}}F\right),
                \ell_2^{m-N} \right)\\
\overset{\refclaim{EigenschaftenDesBanachMazurAbstands}
{EigenschaftenDesBanachMazurAbstands2}}= &\,
d_{BM}\left(\left(F,\restrict{\norm\cdot_{1,m}}F\right),
                    \left(F,\restrict{\norm\cdot_{2,m}}F\right) \right)
 \overset{\eqref{svgeq:5}}\le\sqrt{\frac m k}\le\sqrt{2 t_\eta}\\
\overset{\eqref{svgeq:1}}\le
 &\,c \sqrt{\frac 1 {1-\eps}\log\left(\frac 2 {1-\eps} \right)}
\text.
\end{align*}
\end{proof}